\section{Experimental Evaluation}
\label{sec:experiments}

The results of Section~\ref{sec:analysis} cover \CBSOMminus\ on cliques
with scaled linear biases. In this section, we study through simulation
the settings that those results do not cover: scale-free graphs,
biases outside the region $R$, mixtures of biases, and the memory-based
model \CBSOMplus. We run both models across eight network sizes and a
full range of connectivity levels. Our goals are threefold:
(i)~to test whether only $R$-region bias functions support consensus;
(ii)~to quantify the joint effect of network size and density on the
consensus threshold; and (iii)~to characterize how memory changes these
effects.

% -------------------------------------------------------------------
\subsection{Simulation Setup}
\label{sec:exp-setup}

The simulator generates scale-free graphs by preferential attachment
with degree-distribution exponent $\gamma=2.5$. This value lies in the
empirically documented range for scale-free social
networks~\cite{Barabasi1999,Newman2003,Clauset2009}.
The parameter grid is:
\begin{enumerate}
  \item \textbf{Network sizes:}
        $n \in \{10, 50, 100, 500, 10^3, 10^4, 10^5, 10^6\}$.
  \item \textbf{Density:} $1$--$9$ average neighbors per agent.
  \item \textbf{Bias types:} Unbiased ($\SOMminus$/$\SOMplus$),
        Confirmation, Backfire, and Authority
        (all under $\CBSOMminus$ or $\CBSOMplus$).
  \item \textbf{Runs per configuration:}
        $1{,}000$ ($n \le 10^3$); $500$ ($n=10^4$);
        $50$ ($n=10^5$); $10$ ($n=10^6$).
  \item \textbf{Convergence criterion:}
        opinion spread $\max_i B_i^t - \min_i B_i^t < 10^{-3}$;
        iteration cap $1{,}000$ rounds.
  \item \textbf{Initial opinions:} $B_i^0\sim\mathcal{U}[0,1]$.
\end{enumerate}

\paragraph{Metrics.}
The \emph{Consensus Rate} $\mathrm{CR}$
is the fraction of runs that converge. The
\emph{Mean Convergence Time} $\mathrm{MCT}$ is the average number of
rounds over converging runs. We report $\mathrm{CR}$ throughout.\footnote{All experiments were run on an Intel~i7-1255U machine with 32\,GB RAM and a 512\,GB NVMe SSD.}

% -------------------------------------------------------------------
\subsection{Uniform Networks}
\label{sec:exp-uniform}

A \emph{uniform} network assigns the same bias type to every edge and
the same silence model to every agent. This setting isolates the
contribution of each bias region before we consider mixtures.

\paragraph{Memoryless dynamics ($\CBSOMminus$).}
Figure~\ref{fig:heatmap-uniform} (top row) summarizes the full
$(n,\text{density})$ grid. Both $\SOMminus$ and $\CBSOMminus$
(Confirmation) move from near-zero to near-$100\%$ consensus as
density increases. The Confirmation panel is always at least as light
as the $\SOMminus$ panel: $R$-region biases act as
\emph{consensus accelerators}. The advantage is largest at the
transition. At $n=10^6$ and density~4, the consensus rate is $70\%$ for
$\CBSOMminus$ (Confirmation) and $10\%$ for $\SOMminus$
(Figure~\ref{fig:curves}).
The Backfire panel is uniformly black, i.e., there is no consensus at any
scale or density. Authority also yields zero consensus, except for
negligible events in a 10-agent network. This pattern is consistent with
Theorems~\ref{thm:cbsom-consensus-n3} and~\ref{thm:cbsom-consensus-n2},
which guarantee consensus only for $R$-region biases. However, at
$n=10^6$ the sample is too small (10 runs per configuration) to assert
an absolute dichotomy; these results should be taken as indicative.

The consensus threshold grows sub-linearly with $n$. Density~2 suffices
for $n \le 100$, density~3--4 for $n \in [10^3, 10^5]$, and density~4--5
for $n=10^6$. Figure~\ref{fig:curves} shows the corresponding rightward
shift of the S-curve: as $n$ increases, the transition from zero to full
consensus requires one or two additional neighbors per agent. Once the
threshold is crossed, the consensus rate approaches $100\%$. This holds
in all observed runs for $n \le 10^5$, where each configuration has at
least 50 runs.

\paragraph{Memory-based dynamics ($\CBSOMplus$).}
The bottom row of Figure~\ref{fig:heatmap-uniform} differs markedly
from the top row. For $n \le 100$, the Confirmation panel retains a faint
non-zero region at low densities (e.g., $19.9\%$ at $n=50$,
density~2), while $\SOMplus$ is already near zero.
For $n=500$, $\CBSOMplus$ (Confirmation) peaks at $4.9\%$; for
$n=1{,}000$, it falls to $3.8\%$ and is confined to densities $4$--$7$.
Beyond $n=10^4$, no memory-based configuration, including Confirmation,
reaches consensus in a single run at any density.
Figure~\ref{fig:memory-penalty} makes this collapse visible. The
Confirmation advantage that sustains high consensus in the left panel
($\CBSOMminus$) disappears in the right panel ($\CBSOMplus$)
as $n$ grows.

\begin{figure}[!htbp]
  \centering
  \includegraphics[width=0.85\columnwidth]{fig_heatmap_uniform.pdf}
  \caption{Consensus rate (color, black~$=0\%$, near-white~$=100\%$)
    for \emph{uniform} networks over the full $(n, \text{density})$ grid.
    \textbf{Top row} ($\CBSOMminus$, memoryless): $\SOMminus$ and
    Confirmation transition to near-white at moderate density across all
    eight scales; Backfire and Authority are uniformly black.
    \textbf{Bottom row} ($\CBSOMplus$, memory-based): only Confirmation
    retains a faint non-zero patch for $n \le 100$ and low density;
    all other configurations yield zero}
  \label{fig:heatmap-uniform}
\end{figure}

\begin{figure}[!htbp]
  \centering
  \includegraphics[width=0.85\columnwidth]{fig_consensus_curves.pdf}
  \caption{Consensus rate vs.\ density for $\SOMminus$ (left) and
    $\CBSOMminus$ (Confirmation, right) at six representative scales.
    Confirmation (right) saturates one density unit earlier than $\SOMminus$
    (left) for the same $n$, and the S-curve shifts rightward as
    $n$ grows, reflecting the sub-linear increase in the consensus
    threshold with network size}
  \label{fig:curves}
\end{figure}

\begin{figure}[!htbp]
  \centering
  \includegraphics[width=0.85\columnwidth]{fig_memory_penalty.pdf}
  \caption{Memory penalty under Confirmation bias: $\CBSOMminus$ (left)
    versus $\CBSOMplus$ (right).  In the memoryless model consensus is
    robust across all $n \le 10^3$; in the memory-based model it
    collapses rapidly as $n$ grows, vanishing entirely for $n=1{,}000$
    at high densities}
  \label{fig:memory-penalty}
\end{figure}

% -------------------------------------------------------------------
\FloatBarrier
\subsection{Non-Uniform Networks}
\label{sec:exp-nonuniform}

\emph{Non-uniform} networks assign different bias types to different
edges. We consider four mixed configurations; the proportions refer to
fractions of edges:
\begin{enumerate}
  \item[(i)]  $\SOMminus+\CBSOMminus(\text{Conf.})$, $70\%/30\%$
  \item[(ii)] $\CBSOMminus(\text{Conf.})+\CBSOMminus(\text{Back.})$,
              $80\%/20\%$
  \item[(iii)]$\CBSOMminus(\text{Conf.})+\CBSOMminus(\text{Auth.})$,
              $60\%/40\%$
  \item[(iv)] $\SOMminus+\CBSOMminus(\text{Back.})$, $90\%/10\%$
\end{enumerate}

\paragraph{Rationale for the choice of proportions.}
We selected the four configurations heuristically. They are not
calibrated to measured distributions of cognitive biases in real social
networks; to our knowledge, no such data exist at the resolution
required to assign a bias type to each edge. Instead, the proportions
follow from the uniform-network results of Section~\ref{sec:exp-uniform}
and from an observation already present in the Spiral-of-Silence model
of Aranda et al.~\cite{aranda2024sound}: consensus is structurally
difficult to achieve even without cognitive biases.

Uniform Backfire and Authority configurations yield zero consensus at
every scale and density. Pure non-$R$ configurations therefore offer no
interesting boundary to explore. The relevant question is how small a
disruptive minority suffices to destroy consensus that would otherwise
form. Configuration~(iv) ($10\%$ Backfire) is a \emph{contamination}
scenario: it asks whether a marginal fraction of adversarial edges can
override an otherwise consensus-forming majority. Configuration~(ii)
($20\%$ Backfire) raises the disruptive load but keeps a
receptive-resistant supermajority. It probes the point at which updating
dominated by $\mathrm{conf}(x)$ can no longer absorb the opposing force.
Configuration~(iii) ($40\%$ Authority) places the two regimes in
near-balance. It tests whether substantial hierarchical influence is
compatible with global agreement. Finally, configuration~(i)
($30\%$ $\mathrm{conf}(x)$, $70\%$ unbiased) asks the complementary
question: can a receptive-resistant minority \emph{rescue} or amplify
consensus in an unbiased network that already struggles to converge at
large scales?

The proportions thus span a gradient from marginal perturbation to
near-parity. They cover regimes in which consensus is expected to be
fragile, recoverable, or structurally precluded. Bias distributions
that vary by agent role, community membership, or time may shift these
thresholds. We leave the calibration of proportions to empirical
social-network data for future work.

\paragraph{Memoryless dynamics.}
Figure~\ref{fig:heatmap-nonuniform} shows the results.
Configuration~(i) closely mirrors the uniform Confirmation baseline.
Consensus reaches $97\%$ by density~3 for $n=500$, and the transition
shifts rightward to density~5 only for $n=10^6$. Hence mixing
confirmation-biased and unbiased edges in a $30/70$ proportion
preserves the robustness of consensus.

Configurations (ii)--(iv) collapse to near-zero consensus at every
scale. At $n=100$, configurations (ii) and (iv) produce at most
$2\%$--$3\%$ consensus across all densities; for $n \ge 500$, consensus
is identically zero. A $10\%$ share of Backfire edges
(configuration~iv) suffices to eliminate consensus that would otherwise
occur in almost every run. This \emph{robustness asymmetry} has a clear
direction. Edges in $R$ act as consensus catalysts, whereas edges
outside $R$ prevent consensus, and a surrounding majority of $R$-region
edges cannot neutralize them.

\paragraph{Memory-based dynamics.}
In non-uniform memory-based networks, consensus is essentially absent.
The $\SOMplus+\CBSOMplus(\text{Conf.})$ $70/30$ mixture produces
sporadic consensus only for $n \le 100$ (at most $30\%$, density~4).
For $n \ge 500$, no run converges under any mixed memory-based
configuration. Memory thus amplifies the disruptive effect of non-$R$
edges: at scale, we observed no consensus at all.

\begin{figure}[!htbp]
  \centering
  \includegraphics[width=\columnwidth]{fig_heatmap_nonuniform.pdf}
  \caption{Consensus rate for \emph{non-uniform} memoryless networks
    ($\CBSOMminus$) over the full $(n, \text{density})$ grid.  Only the
    $\SOMminus+\CBSOMminus(\text{Conf.})$ $70/30$ mixture (leftmost panel)
    achieves near-$100\%$ consensus at moderate density. Every
    configuration that contains Backfire or Authority edges collapses to
    near-zero consensus at every scale. This illustrates the robustness
    asymmetry between influence inside and outside the region $R$}
  \label{fig:heatmap-nonuniform}
\end{figure}

% -------------------------------------------------------------------
\FloatBarrier
\subsection{Relative Density}
\label{sec:exp-relative}

The experiments in Sections~\ref{sec:exp-uniform}
and~\ref{sec:exp-nonuniform} measure connectivity by an absolute
neighbor count $k\in\{1,\dots,9\}$. The same count represents very
different levels of social exposure as the network grows.
To assess whether the observed effects depend on this choice, we add a
\emph{relative-density} analysis. In it, the neighborhood of each agent
is a fixed fraction
$\rho\in\{10\%,\,25\%,\,50\%,\,75\%,\,100\%\}$ of the maximum
feasible degree $n-1$.
The $100\%$ condition is the complete-graph limit: every agent
interacts with every other.
Experiments cover $n\in\{100,\,1{,}000,\,10{,}000\}$ with the same
run counts and convergence criterion as Section~\ref{sec:exp-setup}.

Figures~\ref{fig:heatmap-rel-uniform} and~\ref{fig:heatmap-rel-nonuniform}
summarize the four scenarios, one per figure row.

\paragraph{Uniform memoryless (Figure~\ref{fig:heatmap-rel-uniform}, top row).}
$\SOMminus$ and $\CBSOMminus$ (Confirmation) achieve $100\%$ consensus
at every relative density and every tested scale, including the
sparsest configuration ($\rho=10\%$, $n=10{,}000$,
i.e.\ $\approx\!1{,}000$ neighbors per agent).
Backfire and Authority yield $0\%$ everywhere.
The all-white and all-black panels indicate that, in this range, the
bias-region dichotomy depends on the bias type alone and not on the
connectivity level.

\paragraph{Uniform memory-based (Figure~\ref{fig:heatmap-rel-uniform}, bottom row).}
All cells are black for $n\ge 1{,}000$ across all configurations and
all relative densities, including the complete-graph limit
($\rho=100\%$).
The only non-zero values appear at $n=100$: $11.1\%$ and $0.9\%$
for Conf\,+\,Backfire (80/20) at $\rho=10\%$ and $25\%$
respectively, and $0.1\%$ for SOM$^+$\,+\,Confirmation (70/30) at
$\rho=10\%$; all vanish at $\rho\ge 50\%$.

\paragraph{Non-uniform memoryless (Figure~\ref{fig:heatmap-rel-nonuniform}, top row).}
The SOM$^-$\,+\,Conf (70/30) column is entirely white: mixing a $30\%$
confirmation minority with an unbiased majority preserves $100\%$
consensus at every scale and every relative density.
The three disruptive mixtures annotate small non-zero values
($1.4$--$3.0\%$) only at $n=100$; for $n\ge 1{,}000$ they are
identically zero regardless of $\rho$.

\paragraph{Non-uniform memory-based (Figure~\ref{fig:heatmap-rel-nonuniform}, bottom row).}
The entire row is black: zero consensus for all configurations,
all three scales, and all relative densities up to the complete graph.

Together, the two memory-based scenarios sharpen the memory-barrier
conclusion of Section~\ref{sec:exp-discussion}. Increasing connectivity,
up to and including the complete graph, does not overcome the
suppression of consensus under $\CBSOMplus$. The barrier is therefore
structural rather than a threshold effect.

\begin{figure}[!htbp]
  \centering
  \includegraphics[width=0.85\columnwidth]{fig_heatmap_relative_density_uniform.pdf}
  \caption{Consensus rate (color scale as in
    Figures~\ref{fig:heatmap-uniform}--\ref{fig:heatmap-nonuniform})
    for relative-density experiments on \emph{uniform} networks
    ($\rho\in\{10\%,25\%,50\%,75\%,100\%\}$,
    $n\in\{100,\,1\mathrm{k},\,10\mathrm{k}\}$).
    \textbf{Top row} (memoryless): SOM$^-$ and Confirmation are
    near-white at every cell; Backfire and Authority are uniformly
    black; the bias dichotomy does not depend on density.
    \textbf{Bottom row} (memory-based): all black for $n\ge 1{,}000$
    at every $\rho$, including the complete graph ($\rho=100\%$);
    non-zero values appear only at $n=100$, $\rho\le 25\%$
    (11.1\% and 0.9\% for Conf+Backfire; 0.1\% for SOM$^+$+Conf)}
  \label{fig:heatmap-rel-uniform}
\end{figure}

\begin{figure}[!htbp]
  \centering
  \includegraphics[width=0.85\columnwidth]{fig_heatmap_relative_density_nonuniform.pdf}
  \caption{Consensus rate for relative-density experiments on
    \emph{non-uniform} networks (same color scale and parameter
    grid as Figure~\ref{fig:heatmap-rel-uniform}).
    \textbf{Top row} (memoryless): the SOM$^-$+Conf (70/30) column
    is uniformly near-white at all scales and relative densities;
    disruptive mixtures annotate marginal values ($\le 3\%$) only
    at $n=100$, vanishing for $n\ge 1{,}000$.
    \textbf{Bottom row} (memory-based): uniformly black across all
    configurations, scales, and relative densities, including the
    complete-graph limit ($\rho=100\%$)}
  \label{fig:heatmap-rel-nonuniform}
\end{figure}

% -------------------------------------------------------------------
\FloatBarrier
\subsection{Discussion}
\label{sec:exp-discussion}

Three empirical findings summarize the experimental evidence:

\begin{enumerate}
  \item \textbf{Bias-region dichotomy.}
    $R$-region biases (Confirmation) yield $\mathrm{CR}$ approaching
    $100\%$ at sufficient density for $n\le 10^5$ (at least 50 runs per
    configuration). At $n=10^6$ (10~runs), all Confirmation runs converge
    at density~$\ge 5$, but the sample is too small to assert
    $\mathrm{CR}=100\%$ unconditionally.
    Non-$R$ biases (Backfire, Authority) yield $\mathrm{CR}=0\%$
    in all observed runs at every scale and density tested.

  \item \textbf{Sub-linear consensus threshold.}
    The critical density grows from $2$ (small networks) to $4$--$5$
    ($n=10^6$), spanning only $3$ units across eight orders of
    magnitude of network size. The threshold under Confirmation is
    consistently one unit lower than under the unbiased $\SOMminus$
    baseline.

  \item \textbf{Memory as a structural barrier.}
    Memory-based dynamics ($\SOMplus$, $\CBSOMplus$) almost entirely
    suppress consensus. Even Confirmation retains non-zero
    $\mathrm{CR}$ only at $n \le 1{,}000$. In non-uniform
    settings, an $R$-region majority cannot compensate for the path
    dependence induced by memory beyond $n=100$.
\end{enumerate}

\paragraph{Data and code availability.}
The Scala/Akka simulator used to produce the results reported in this
section is publicly available at
\url{https://github.com/juanmarcosdev/belief_evolution_simulator}.
